\section{The Bulk}

The curved four-dimensional interior of the mesh is the bulk. It is hyperbolic, which means it curves so that there is more and more room the farther out you go, like a tree branching faster than a flat grid ever could. That extra room does real work in what follows.

\begin{figure}[ht]
\centering
\begin{minipage}{0.21\textwidth}
\centering
\includegraphics[width=\linewidth]{curvature-spherical}
\\[2pt]{\small positive, closed}
\end{minipage}\hspace{0.03\textwidth}
\begin{minipage}{0.21\textwidth}
\centering
\includegraphics[width=\linewidth]{curvature-euclidean}
\\[2pt]{\small zero, flat}
\end{minipage}\hspace{0.03\textwidth}
\begin{minipage}{0.21\textwidth}
\centering

\includegraphics[width=\linewidth]{curvature-hyperbolic}
\\[2pt]{\small negative, open}
\end{minipage}
\caption{The three constant curvatures. The base is hyperbolic, the open case on the right, where space spreads faster than it does when flat.}
\label{fig:curvatures}
\end{figure}

Because the interior swells exponentially while a region's surface grows only as its area, almost everything in a region sits near its edge, and information about the interior is crowded outward onto the boundary with nowhere else to go. That crowding is the holographic principle, here a plain geometric fact rather than a deep mystery. It is also why a bounded pattern in the bulk can hold far more inside than its size in flat space would ever allow, the room that a mind or a memory needs.

The bulk has a definite four-dimensional size, which is forced rather than measured by a naive ruler. A strongly curved honeycomb has no flat regime, its balls grow exponentially at every scale, so the usual dimension estimators are confounded by the curvature and read too high. The right way to read the dimension is on the flat boundary instead (the next section), where a Euclidean regime exists, and the bulk is then one dimension higher because the boundary is a slice of it. Read that way, the bulk is four-dimensional and the boundary three, exactly as the rank of $D_4$ demands. The curved bulk is also where the radial direction is not a place but a scale: going deeper is the same as standing back and squinting, the holographic scale axis, which is why one cannot coarse-grain along it without blurring across scales.
